\section{Related Work}\label{sec:related-work}

Graph neural networks can support intrusion detection by modeling relationships among network and system entities. Pujol-Perich et al.~\cite{pujolperich2021robustgnnids} show that graph structure can make network intrusion detection more robust than models that classify flows independently. E-GraphSAGE~\cite{lo2022egraphsage} applies a GNN to flow-based intrusion detection in Internet of Things environments and uses both edge features and graph structure. Recent surveys show that graph-based intrusion detection is a growing area, but that deployment, labeling, and robustness remain open challenges~\cite{GNN_IDS}. These works support the use of graphs for security monitoring. However, many graph-based intrusion detectors are supervised and depend on labeled non-malicious and malicious examples. In contrast, \sysname learns representations from non-malicious graphs.

Provenance and temporal graph methods are also closely related to our task. These methods use graph structure to describe how system activity changes over time. Jbeil~\cite{Jbeil} studies temporal graph learning for lateral movement detection in enterprise networks. MAGIC~\cite{MAGIC} uses masked graph representation learning to detect advanced persistent threats from provenance graphs. FLASH~\cite{Flash} and KAIROS~\cite{cheng2024kairos} also use provenance graphs to support intrusion detection and investigation. These systems show that graph-structured representations are useful for security monitoring. Our work differs in focus. We study graph-level anomaly scoring for benign-only representation learning and focus on how topology-aware contrastive embeddings should be used for detection.

Graph contrastive learning is significant for this problem because it can learn useful graph embeddings without attack labels. GraphCL~\cite{you2020graphcl} learns graph-level representations by creating augmented views of the same graph and training an encoder to make their embeddings agree. This makes GraphCL an important baseline for our work. However, GraphCL does not explicitly model higher-order graph shape. This can be a limitation when non-malicious and malicious activity graphs share the same broad structure but differ in smaller structural details.

TopoGCL~\cite{chen2024topogcl} addresses this limitation by adding topological information to graph contrastive learning. It uses topological summaries to describe graph shape across multiple scales. This makes TopoGCL the closest methodological starting point to \sysname. However, TopoGCL was designed mainly for unsupervised graph classification, not anomaly scoring or intrusion detection. A graph classifier can use labeled examples after representation learning to separate classes. A model trained only on non-malicious graphs cannot rely on labeled attacks during representation learning. It must decide whether a test graph is abnormal by comparing it with non-malicious training graphs.

\sysname builds on graph contrastive learning and topology-aware graph learning, but it targets a different security setting. It adapts topology-aware contrastive learning for intrusion detection and scores test graphs by their distance from non-malicious training embeddings. The key difference is that \sysname scores graphs in the same projected embedding space shaped by the contrastive loss. This keeps representation learning and anomaly scoring consistent. The goal is not to introduce a new topological summary. Instead, \sysname uses topology-aware contrastive embeddings more directly for intrusion detection in intelligent systems.
